% !TeX root=../main.tex
\chapter{جمع‌بندی و پیشنهادهای آینده}
\label{chap:conclusion}

\section{جمع‌بندی تحقیق}
این رساله مسئلهٔ شناسایی و بازسازی کور ساختار کدهای توربو در مخابرات غیرهمکارانه را صورت‌بندی کرد و روشی برای بازسازی مشترک پلی‌نوم‌های RSC و درهم‌ریزهای داخلی و خارجی با لحاظ خاتمهٔ ترلیس و نویز ارائه داد.

\section{دستاوردهای اصلی}
دستاوردها شامل چارچوب ریاضی پارامترهای مجهول، الگوریتم جست‌وجوی درختی مقید دوفازی، معیارهای آماری هرس، تحلیل ابهام ساختاری، و ارزیابی نظری/تجربی پیچیدگی و عملکرد است.

\section{پاسخ به پرسش‌های تحقیق}
در این بخش، پاسخ‌های مستند به پرسش‌های مطرح‌شده در فصل اول بر اساس نتایج فصول چهارم تا ششم جمع‌بندی می‌شود.

\section{محدودیت‌های روش}
محدودیت‌هایی مانند حساسیت به نویز بالا، هزینهٔ محاسباتی در طول‌های بزرگ، وابستگی به فرضیات مدل مشاهده، و ابهام ساختاری باقی‌مانده بیان می‌گردد.

\section{کاربردهای عملی}
دلالت‌های روش برای سامانه‌های غیرهمکارانه، پایش طیف و مهندسی معکوس کد کانال مرور می‌شود.

\section{مسیرهای پژوهشی آینده}
پیشنهادهای آینده شامل تعمیم به کدهای مرتبط، بازسازی برخط، بهره‌گیری قوی‌تر از اطلاعات نرم، یادگیری برای هرس، و ارزیابی روی داده‌های واقعی میدان است.
